\documentclass[12pt]{ximera}
\title{Final Exam}
\begin{document}
\begin{abstract}
Final exam, covering apportionment and game theory: how quotas respond to a changing
divisor, which paradoxes each method can suffer, a strictly determined game, a
four-neighborhood apportionment by Jefferson, Adams, and Huntington-Hill, and a
two-finger game at equilibrium --- plus bonus review from earlier chapters.
\end{abstract}
\maketitle

\section*{Final Exam}

Round all decimals to two places unless otherwise stated. Useful formulas:
\[
SD=\frac{P}{M}, \qquad q_i=\frac{P_i}{SD}, \qquad \text{geometric mean of } a,b=\sqrt{ab},
\]
\[
E(A_1)=P(B_1)\cdot(\text{value of }A_1B_1)+P(B_2)\cdot(\text{value of }A_1B_2).
\]

\begin{problem}
Assess whether each of the following claims is true or false.

\begin{prompt}
\textbf{(a)} Claim: the quotas of states with the largest populations change the fastest
when the standard divisor of an apportionment problem is decreased.

\begin{multipleChoice}
\choice[correct]{True}
\choice{False}
\end{multipleChoice}

\begin{feedback}
True. A quota is $P_i/d$, so changing the divisor scales every quota by the same factor.
The same percentage change is a larger absolute change for a larger quota --- which
belongs to a larger state.
\end{feedback}
\end{prompt}

\begin{prompt}
\textbf{(b)} Claim: the quotas of states with the smallest populations change the fastest
when the standard divisor of an apportionment problem is increased.

\begin{multipleChoice}
\choice{True}
\choice[correct]{False}
\end{multipleChoice}

\begin{feedback}
False. Whether the divisor goes up or down, every quota is scaled by the same factor, so
the largest states' quotas always move the most.
\end{feedback}
\end{prompt}

\begin{prompt}
\textbf{(c)} Claim: every game that can be represented with a $2\times 2$ payoff matrix
must have a saddle point, an outcome where the players' dominant strategies intersect.

\begin{multipleChoice}
\choice{True}
\choice[correct]{False}
\end{multipleChoice}

\begin{feedback}
False. Only strictly determined games have a saddle point. The finger game in
Problem 5 has none, so both players must mix their strategies.
\end{feedback}
\end{prompt}

\begin{prompt}
\textbf{(d)} Claim: every game that can be represented with a $2\times 2$ payoff matrix
must have an equilibrium point, where neither player has reason to divert from their
chosen strategy.

\begin{multipleChoice}
\choice[correct]{True}
\choice{False}
\end{multipleChoice}

\begin{feedback}
True. If there is no saddle point, the players' optimal \emph{mixed} strategies form an
equilibrium: neither can do better by changing their proportions alone.
\end{feedback}
\end{prompt}
\end{problem}

\begin{problem}
Determine all applicable answers to the following.

\begin{prompt}
\textbf{(a)} Which of the following is possible in an apportionment using Jefferson's
method?

\begin{selectAll}
\choice[correct]{Upper quota violation}
\choice{Lower quota violation}
\choice{Population paradox}
\choice{New States paradox}
\end{selectAll}

\begin{feedback}
Jefferson's method rounds modified quotas down with a divisor \emph{smaller} than the
standard one, which can push a large state above its upper quota --- but never below its
lower quota. Like every divisor method, it is free of the Alabama, Population, and New
States paradoxes.
\end{feedback}
\end{prompt}

\begin{prompt}
\textbf{(b)} Which of the following is possible in an apportionment using Hamilton's
method?

\begin{selectAll}
\choice{Upper quota violation}
\choice{Lower quota violation}
\choice[correct]{Alabama paradox}
\choice[correct]{New States paradox}
\end{selectAll}

\begin{feedback}
Hamilton's method always satisfies the quota rule, since each state gets its lower quota
plus at most one extra seat. The price is that it can suffer all three paradoxes,
including the Alabama and New States paradoxes.
\end{feedback}
\end{prompt}
\end{problem}

\begin{problem}
Consider the game associated with the following payoff matrix.
\[
\begin{bmatrix} 2 & -7\\ 5 & 3 \end{bmatrix}
\]

\begin{prompt}
\textbf{(a)} Claim: this is a strictly determined game.

\begin{multipleChoice}
\choice[correct]{True}
\choice{False}
\end{multipleChoice}

\begin{feedback}
True. The row minimums are $-7$ and $3$, with maximum $3$; the column maximums are $5$
and $3$, with minimum $3$. These agree, so the entry $3$ (row 2, column 2) is a saddle
point.
\end{feedback}
\end{prompt}

\begin{prompt}
\textbf{(b)} What is the value of this game?

\begin{multipleChoice}
\choice{$5$}
\choice[correct]{$3$}
\choice{$2$}
\choice{$-7$}
\choice{None of these}
\end{multipleChoice}

\begin{feedback}
The value of a strictly determined game is its saddle point entry: $3$.
\end{feedback}
\end{prompt}
\end{problem}

\begin{problem}
A class of $110$ architecture students is studying a small city with four main
neighborhoods. Each student will be assigned to one neighborhood to construct part of a
scale model of the entire city.

\begin{center}
\begin{tabular}{|r|c|c|c|c|c|}\hline
\textbf{Neighborhood} & A & B & C & D & \textbf{Total}\\\hline
\textbf{Population} & 9,455 & 752 & 245 & 548 & 11,000\\\hline
\end{tabular}
\end{center}

\begin{prompt}
\textbf{(a)} Compute the standard divisor and standard quotas.

Standard divisor: $\answer{100}$

\smallskip
$q_A=\answer[tolerance=0.005]{94.55}$ \quad $q_B=\answer[tolerance=0.005]{7.52}$ \quad
$q_C=\answer[tolerance=0.005]{2.45}$ \quad $q_D=\answer[tolerance=0.005]{5.48}$

\begin{feedback}
\[ SD=\frac{11{,}000}{110}=100, \]
so each quota is the population divided by $100$.
\end{feedback}
\end{prompt}

\begin{prompt}
\textbf{(b)} Round the standard quotas as if doing just the first step of each method.

Jefferson (down): $A: \answer{94}$, $B: \answer{7}$, $C: \answer{2}$, $D: \answer{5}$,
total $\answer{108}$

\smallskip
Adams (up): $A: \answer{95}$, $B: \answer{8}$, $C: \answer{3}$, $D: \answer{6}$,
total $\answer{112}$

\smallskip
Huntington-Hill: $A: \answer{95}$, $B: \answer{8}$, $C: \answer{3}$, $D: \answer{6}$,
total $\answer{112}$

\begin{feedback}
Jefferson is $2$ students short, and Adams is $2$ over.

For Huntington-Hill, compare each quota with the geometric mean of its neighbors:
\[
\sqrt{94\cdot95}\approx 94.4987,\quad \sqrt{7\cdot8}\approx 7.4833,\quad
\sqrt{2\cdot3}\approx 2.4495,\quad \sqrt{5\cdot6}\approx 5.4772.
\]
All four quotas lie just above their cutoffs ($94.55$, $7.52$, $2.45$, $5.48$), so all
round up --- $2$ students over.
\end{feedback}
\end{prompt}

\begin{prompt}
\textbf{(c)} Complete the apportionment by each method, choosing the appropriate modified
divisor from $98$, $100.053$, and $102$.

\textbf{Jefferson} --- divisor $\answer{98}$, giving $A: \answer{96}$, $B: \answer{7}$,
$C: \answer{2}$, $D: \answer{5}$

\smallskip
\textbf{Adams} --- divisor $\answer{102}$, giving $A: \answer{93}$, $B: \answer{8}$,
$C: \answer{3}$, $D: \answer{6}$

\smallskip
\textbf{Huntington-Hill} --- divisor $\answer[tolerance=0.0005]{100.053}$, giving
$A: \answer{95}$, $B: \answer{8}$, $C: \answer{2}$, $D: \answer{5}$

\begin{hint}
Jefferson is short and needs a smaller divisor; Adams and Huntington-Hill are over and
need larger ones.
\end{hint}

\begin{feedback}
\textbf{Jefferson}, $d=98$: modified quotas $96.48$, $7.67$, $2.50$, $5.59$; rounding
down gives $96+7+2+5=110$. \checkmark

\textbf{Adams}, $d=102$: modified quotas $92.70$, $7.37$, $2.40$, $5.37$; rounding up
gives $93+8+3+6=110$. \checkmark

\textbf{Huntington-Hill}, $d=100.053$ (four decimals):
\[
\begin{array}{lcccc}
 & A & B & C & D\\
mq & 94.4999 & 7.5160 & 2.4487 & 5.4771\\
\text{cutoff} & 94.4987 & 7.4833 & 2.4495 & 5.4772\\
\text{rounded} & 95 & 8 & 2 & 5
\end{array}
\]
for a total of $110$. \checkmark\ Neighborhoods C and D fall just below their cutoffs,
by less than $0.001$.

Notice that Jefferson gives the large neighborhood A $96$ students, above its upper
quota of $95$ --- an upper quota violation, as in Problem 2(a).
\end{feedback}
\end{prompt}
\end{problem}

\begin{problem}
Ahmad and Bojan are playing a game where on the count of three they each hold up one or
two fingers. If the sum is even, Ahmad wins points equal to the total; if the sum is
odd, Bojan wins points equal to the total. The payoff matrix (to Ahmad) is
\[
\begin{array}{r|rr}
 & \text{B: one} & \text{B: two}\\\hline
\text{A: one} & 2 & -3\\
\text{A: two} & -3 & 4
\end{array}
\]

\begin{prompt}
\textbf{(a)} What proportion of the time will Ahmad show one finger at the equilibrium
point? $\answer{7/12}$

\begin{hint}
Choose Ahmad's proportion $p$ of ``one'' so that Bojan's two options give Ahmad the same
expected payoff.
\end{hint}

\begin{feedback}
Against Bojan's ``one'': $2p-3(1-p)=5p-3$. Against Bojan's ``two'':
$-3p+4(1-p)=4-7p$. Setting these equal, $12p=7$, so $p=\frac7{12}$.
\end{feedback}
\end{prompt}

\begin{prompt}
\textbf{(b)} What proportion of the time will Ahmad show two fingers?
$\answer{5/12}$

\begin{feedback}
$1-\frac7{12}=\frac5{12}$.
\end{feedback}
\end{prompt}

\begin{prompt}
\textbf{(c)} What proportion of the time will Bojan show one finger at the equilibrium
point? (The symmetry of the matrix should save you some work.) $\answer{7/12}$

\begin{feedback}
The matrix is symmetric, so Bojan's equation is the same as Ahmad's: with proportion $q$
of ``one'', $2q-3(1-q)=-3q+4(1-q)$ gives $q=\frac7{12}$.
\end{feedback}
\end{prompt}

\begin{prompt}
\textbf{(d)} What proportion of the time will Bojan show two fingers?
$\answer{5/12}$

\begin{feedback}
$1-\frac7{12}=\frac5{12}$.
\end{feedback}
\end{prompt}

\begin{prompt}
\textbf{(e)} If both players play optimally, what is the expected value of this game?
$\answer{-1/12}$

\begin{feedback}
Using Ahmad's strategy against either of Bojan's options:
\[
5\cdot\tfrac7{12}-3=\tfrac{35-36}{12}=-\tfrac1{12}.
\]
The game slightly favors Bojan: about $0.08$ points per round.
\end{feedback}
\end{prompt}
\end{problem}

\begin{problem}
The following are bonus questions from earlier chapters.

\begin{prompt}
\textbf{(a)} Identify the Condorcet candidate in the following election.
\begin{center}
\begin{tabular}{c|ccc}
Number of votes & 8 & 6 & 4\\\hline
1st & B & A & C\\
2nd & C & D & B\\
3rd & A & C & D\\
4th & D & B & A
\end{tabular}
\end{center}

\begin{multipleChoice}
\choice{A}
\choice{B}
\choice[correct]{C}
\choice{D}
\choice{There is none}
\end{multipleChoice}

\begin{feedback}
C beats B $10$--$8$, A $12$--$6$, and D $12$--$6$.
\end{feedback}
\end{prompt}

\begin{prompt}
\textbf{(b)} In a weighted voting system with $11$ players, how many coalitions include
$P_1$, $P_2$, and $P_3$ but exclude $P_{10}$ and $P_{11}$? $\answer{64}$

\begin{feedback}
Five players are settled, leaving $6$ free: $2^6=64$.
\end{feedback}
\end{prompt}

\begin{prompt}
\textbf{(c)} Consider the weighted voting system $[10:7,5,3,1]$ with players A, B, C, and
D. Which are the winning coalitions?

\begin{selectAll}
\choice[correct]{$\set{A,B}$}
\choice[correct]{$\set{A,C}$}
\choice{$\set{A,D}$}
\choice{$\set{B,C,D}$}
\choice[correct]{$\set{A,B,C}$}
\choice[correct]{$\set{A,B,D}$}
\choice[correct]{$\set{A,C,D}$}
\choice[correct]{$\set{A,B,C,D}$}
\end{selectAll}

\begin{feedback}
The weights are $\set{A,B}=12$, $\set{A,C}=10$, $\set{A,B,C}=15$, $\set{A,B,D}=13$,
$\set{A,C,D}=11$, $\set{A,B,C,D}=16$. $\set{A,D}=8$ and $\set{B,C,D}=9$ fall short.
\end{feedback}
\end{prompt}

\begin{prompt}
\textbf{(d)} Bandit, Chili, Bingo, and Bluey are having a picnic. Chili splits the food
into four shares she views as equal, and the others view only the following shares as
fair:
\begin{center}
\begin{tabular}{c|ccc}
Name & Bandit & Bingo & Bluey\\\hline
Bids & $\set{s_1,s_2,s_3}$ & $\set{s_2,s_3}$ & $\set{s_1,s_3}$
\end{tabular}
\end{center}
List three different ways of achieving a fair division.

\begin{freeResponse}
\end{freeResponse}

\begin{feedback}
Nobody bids on $s_4$, so Chili, the divider, takes it. Then:
\begin{itemize}
\item Bingo $s_2$, Bluey $s_1$, Bandit $s_3$;
\item Bingo $s_2$, Bluey $s_3$, Bandit $s_1$;
\item Bingo $s_3$, Bluey $s_1$, Bandit $s_2$.
\end{itemize}
These are the only three.
\end{feedback}
\end{prompt}

\begin{prompt}
\textbf{(e)} Give a two-sentence summary of Hamilton's method of apportionment,
including what specifically distinguishes it from the other methods studied.

\begin{freeResponse}
\end{freeResponse}

\begin{feedback}
For example: Hamilton's method gives each state its lower quota, then hands the
leftover seats one at a time to the states with the largest fractional parts. Unlike
the divisor methods (Jefferson, Adams, Webster, Huntington-Hill), it never changes the
divisor, so it always satisfies the quota rule --- but it is vulnerable to the Alabama,
Population, and New States paradoxes.
\end{feedback}
\end{prompt}
\end{problem}

\end{document}
